\documentclass[exercice,a4]{thedocument}%
\usepackage[french]{babel}%
\usepackage{tabularx}%
\usepackage{array}%
\usepackage[np]{numprint}%
\usepackage{amsmath}%
\startdatakeys
    \source{22GENMATNC1}
    \BO{2026}
    \cycle{4}
    \competencesDuSocle{CA,RA}
    \objectifsApprentissage{C4C2aa,C4C2ag,C4B3bd}
    \motsCles{tableur,moyenne,coordonnées géographiques,volume,boule}
\enddatakeys
\begin{document}%

\medskip

Cet exercice est un questionnaire à choix multiples (QCM).

Pour chaque question, \textbf{une seule} des trois réponses proposées est exacte.

\textbf{Sur la copie, indiquer le numéro de la question et la réponse A, B ou C choisie.}

\textbf{Aucune justification n'est demandée.}

Aucun point ne sera enlevé en cas de mauvaise réponse.

\begin{center}
\begin{tabularx}{\linewidth}{|c|m{5cm}|*{3}{>{\centering \arraybackslash}X|}}\hline
\multicolumn{2}{|c|}{Questions}&Réponse A&Réponse B&Réponse C\\ \hline
1&Dans un tableur, quelle formule faut-il saisir dans la cellule D1 pour afficher la somme des nombres des cellules A1, B1 et C1 ?

\begin{tabular}{|*{5}{c|}}\hline
	&A&B&C&D\\ \hline
1	&3&5&4&\\ \hline
\end{tabular}&=somme(A1:~C1)&=(A1:~C1)&somme(A1*C1)\\ \hline
2&Soit la série de nombres:

15 ; 10 ; 13 ; 9 ; 10 ; $x$.

La moyenne de la série est 11 pour $x$
égal à \ldots&9&10&11\\ \hline
3&Sur la Terre, l'équateur est:&un méridien&un demi-cercle&un parallèle\\ \hline
4&Le volume exact, en cm$^3$, d'une boule de $6$~cm de diamètre est :

{\footnotesize On rappelle le volume $V$ d'une boule de rayon $R$
: $V = \dfrac{4\pi R^3}{3}$}\rule[-3mm]{0mm}{9mm}&$36\pi$&\np{113,0973355}&$288\pi$\\ \hline
\end{tabularx}
\end{center}

\bigskip

\showthesource
\end{document}
